\section{The Spitzer--Cox example of phase transition}

\begin{definition}[Georgii (11.20)--(11.27)]
    \label{def:sc-matrix}
    \uses{def:bl-transfer-matrix}
    \lean{MeasureTheory.GibbsMeasure.Markov.SpitzerCox.binomialWeight, MeasureTheory.GibbsMeasure.Markov.SpitzerCox.poissonWeight, MeasureTheory.GibbsMeasure.Markov.SpitzerCox.sum_binomialWeight_mul_binomialWeight, MeasureTheory.GibbsMeasure.Markov.SpitzerCox.hasSum_binomialWeight_mul_binomialWeight, MeasureTheory.GibbsMeasure.Markov.SpitzerCox.matrix, MeasureTheory.GibbsMeasure.Markov.SpitzerCox.isTransferMatrix}
    \leanok

    On $E = \mathbb Z_+$ the branching-with-immigration matrix
    $Q(x,y) = \sum_k \ell(x,p,k)\mathfrak p(1-p, y-k)$ is a transfer matrix. The binomial weights
    convolve, $\sum_{j \leq k}\ell(m,p,j)\ell(n,p,k-j) = \ell(m+n,p,k)$ (11.22), and compose,
    $\sum_k \ell(n,p_1,k)\ell(k,p_2,j) = \ell(n,p_1p_2,j)$ (11.23).
\end{definition}

\begin{proposition}[Georgii (11.28)]
    \label{prop:sc-11-28}
    \uses{def:sc-matrix}
    \lean{MeasureTheory.GibbsMeasure.Markov.SpitzerCox.poissonWeight_one_mul_matrixReal_comm}
    \leanok

    $Q$ is reversible with respect to the unit Poisson weights: $\alpha(x)Q(x,y) = \alpha(y)Q(y,x)$
    with $\alpha = \mathfrak p(1,\cdot)$. No condition on $p$ is needed.
\end{proposition}

\begin{theorem}[Georgii (11.29)--(11.30)]
    \label{thm:sc-11-29}
    \uses{prop:sc-11-28, def:bl-boundary-law, thm:bl-11-9a}
    \lean{MeasureTheory.GibbsMeasure.Markov.SpitzerCox.rightReal, MeasureTheory.GibbsMeasure.Markov.SpitzerCox.isBoundaryLaw, MeasureTheory.GibbsMeasure.Markov.SpitzerCox.chain, MeasureTheory.GibbsMeasure.Markov.SpitzerCox.chain_intervalCylinder_self, MeasureTheory.GibbsMeasure.Markov.SpitzerCox.eq_of_chain_eq, MeasureTheory.GibbsMeasure.Markov.SpitzerCox.not_countable_G}
    \leanok

    For $u, v \geq 0$ the pair $\ell^u_i, r^v_i$ is a boundary law, whose Gibbs measure satisfies
    $\mu^{u,v}(\sigma_i = x) = \mathfrak p((1+up^i)(1+vp^{-i}), x)$. Distinct parameters give
    distinct measures, so $\mathcal G(\gamma^Q)$ is uncountable.
\end{theorem}
\begin{proof}
    \leanok
    $Qr^v_i = r^v_{i-1}$ and $\ell^u_i r^v_i = e^{uv}$ by the binomial identities; the one-site
    marginal (11.30) at three sites separates the parameters.
\end{proof}

\begin{lemma}[The Poisson limit theorem, as used in (11.31) and (11.46)]
    \label{lem:sc-poisson-limit}
    \lean{ProbabilityTheory.tendsto_choose_mul_pow_of_tendsto_mul, ProbabilityTheory.tendsto_binomial_real_singleton_of_tendsto_mul, ProbabilityTheory.tendsto_tsum_abs_binomial_real_sub_of_tendsto_mul, ProbabilityTheory.tendsto_binomial_of_tendsto_mul, tendsto_finsetSum_of_dominated_convergence, ENNReal.tendsto_sum_range_of_dominated_convergence}
    \leanok

    If $x_n p_n \to c$ and $p_n \to 0$ along any filter, then $\ell(x_n, p_n, k) \to \mathfrak p(c,k)$
    for every $k$, and the binomial laws converge to the Poisson law in total variation. Along a
    growing family of finite sums, dominated convergence holds with a summable bound.
\end{lemma}
\begin{proof}
    \leanok
    $x\log(1+g) - xg$ is $O(|xg|\,|g|)$, and the falling factorial $\prod_{i<k}(x_n - i)p_n \to c^k$;
    Scheff\'e's lemma for series upgrades pointwise to total-variation convergence. Mathlib's
    version fixes the number of trials to the index and passes through an asymptotic equivalence
    that fails when the trial count stays bounded, so it does not transfer.
\end{proof}
